\documentclass[../main.tex]{subfiles}

\begin{document}

\chapter{Conclusion}

Well that was a journey. In 9 weeks, I have learned how to build a fully functioning 
(but definitely not industry-level) SAT solver, and hopefully in the amount of time it 
took you to read through all of this, you did too. Despite how much we've done, there is 
a much larger world of SAT solving out there. CDCL and DPLL are just two algorithms out 
of many that can be used to solve SAT problems. Some are better than others in certain 
cases. That is why, like some ML models, SAT solvers also use an ``ensemble-like'' method 
of solving. Unlike ML ensembles, which takes the output of multiple models and ``weighs''
them to form the final output, these SAT solvers run multiple different solving 
algorithms and heuristics at once, and whichever gets an answer first is the winner. 
This is because those algorithms will all have the same output, some just get there 
faster.

There are also parallel algorithms for SAT solving, so instead of one thread solving 
one problem, you can have multiple threads solving one problem or multiple threads solving 
multiple problems. 

Outside of solving, there is also the notion of simplification of the problem in general. 
These preprocessing and inprocessing algorithms are capable of removing redundant clauses 
and variables from the problem entirely, thus making it easier to solve. Likewise, we can 
also strengthen clauses. Even cooler, we can solve a CNF problem, add more clauses and variables,
then use our previous solution to try to help solve those new constraints. This is called 
incremental SAT, which is a really important feature.

Obviously if I had more than 9 weeks
\footnote{It was a self-imposed timeline, but that's irrelevant.}
to write and learn about SAT solvers, we could learn how to extend it even further, 
perhaps even one day being competition worthy. There is so much more out there, and I 
think it would be amazing if we got to explore that more in a second draft.

\section{Improvements}

Before I end this for good, I want to list some improvements needed for this paper
and our solving.

\subsection{This Paper is Poorly Written}

You know it, I know it, everyone knows it. This would never pass academic review. 
You're thinking ``this guy has never written anything before'' and you're 100\% 
right, I haven't. I need to go back, review it, reformat, add actual citations, etc.
Is it my fault it's in a poorly written state to begin with? 100\%, and I while I won't 
be making this project my entire life until I die, I will go back every now and then to 
fix things, reword stuff, reformat, whatever needs done. You can help me out by first 
reading this, then contacting me on revisions I should make. Any advice and help is 
really appreciated, and I'm writing this for the people who are interested, so 
by showing interest in this, it tells me how much more I need to put in to make it 
better for y'all.

One thing I'd like to do is redo the chapter on CDCL so it's not one mega-chapter
and it could be separated into its necessary parts (clause learning, backjumping)
and then the things to improve it those aspects (VSIDS, first UIPs, restarts)
could be in their own chapter, then lastly, a chapter on the things that we 
need to keep in mind when we implement on actual hardware (clause database management).
This would make the flow of everything more isolated and easier to follow along 
rather than following in my exact footsteps. 

\subsection{The Code Needs to be Cleaned}

I said to myself ``I'll refactor this nicely later'' and later never came, and before
I knew it, I've finished writing everything I set out to write! The code is supposed to 
be used as a teaching aide, and when it's not very clean, it's not as good as it could be. 
Sure there are code comments to help explain some things where the implementation differ 
from the pseudocode, but that doesn't make the code more readable, just easier to understand 
if you \textit{can} read it.

Speaking of pseudocode, it's not in great shape. The pseudocode is really similar to the 
actual code, which is a strength and a curse. I'm not great at writing pseudocode, and I 
think that it's a skill that I need to develop more. If I were better at writing pseudocode,
the first thing I would do is reformat it, then make it less C-based so it's actual 
pseudocode instead of ``slightly more abstract C++.'' This would tie in better to the 
theory of SAT solving, then the actual code could be used as a teaching aide on how it's 
implemented.

\subsection{We Don't Explore Other Real Solvers}

Since this is a paper on how to build a solver, the main focus was on how to build one, 
but we ignore how other solvers have been built. This could be a good chance to explore 
the history of SAT solving, from the people, the algorithms, and the programs that have 
been built throughout the years. And through exploring other solvers and algorithms, we 
could have new chapters and sections on how they work and how to implement them, which 
would be pretty cool.

\subsection{Our Solver is C++ Only}

Our SAT solver is C++ only. It will only run in C++, and depending on how you write it, 
it could be limited to just C++20, if you don't precompile it. This means if we wanted 
to use it in other languages like C, Java, JavaScript, Python, Rust, or any other language, 
we would need to create a bridge between the two languages. This isn't a very difficult 
undertaking, however, as the solver grows, we would need to update the bridge for each 
language, which wouldn't be fun. Since we also deal with implementation details in this 
book, building bridges between languages could be a fun thing to write and learn about.

\subsection{There Are No Actual Citations}

This one is pretty self explanatory. Even though I tell you where I got my facts and 
figures, I don't properly cite them. I'm pretty sure this paper would be enough to 
get me expelled for plagiarism due to how I don't use proper citations anywhere. 

And with that, those are all I can think of that needs to be done for this paper. 
Thank you for reading until the end. Til next time or something.

\end{document}